\documentclass[conference]{IEEEtran}
\IEEEoverridecommandlockouts
\usepackage{cite}
\usepackage{amsmath,amssymb,amsfonts}
\usepackage{graphicx} 
\usepackage{textcomp}
\usepackage{xcolor}
\usepackage{booktabs}
\usepackage{array}
\usepackage{tabulary}
\usepackage{ragged2e}
\usepackage{makecell}
\usepackage{balance}
\usepackage{url}
\usepackage{flushend}
\usepackage{placeins}
\def\BibTeX{{\rm B\kern-.05em{\sc i\kern-.025em b}\kern-.08em
    T\kern-.1667em\lower.7ex\hbox{E}\kern-.125emX}}
\begin{document}
\title{A Feature-Reduced SSH Brute-Force Detection Method Using XGBoost on AWS Infrastructure}

\author{\IEEEauthorblockN{1\textsuperscript{st} Hero Yudo Martono}
\IEEEauthorblockA{\textit{Department of Informatics and Computer Engineering} \\
\textit{Politeknik Elektronika Negeri Surabaya}\\
Surabaya, Indonesia \\
hero@pens.ac.id}
\and
\IEEEauthorblockN{2\textsuperscript{nd} Reesa Akbar}
\IEEEauthorblockA{\textit{Electrical Engineering} \\
\textit{Politeknik Elektronika Negeri Surabaya}\\
Surabaya, Indonesia \\
reesa@pens.ac.id}
\and
\IEEEauthorblockN{3\textsuperscript{rd} Dian Septiani Santoso}
\IEEEauthorblockA{\textit{Department of Informatics and Computer Engineering} \\
\textit{Politeknik Elektronika Negeri Surabaya}\\
Surabaya, Indonesia \\
dian@pens.ac.id}
\and
\IEEEauthorblockN{4\textsuperscript{th} Andhik Ampuh Yunanto}
\IEEEauthorblockA{\textit{Department of Informatics and Computer Engineering} \\
\textit{Politeknik Elektronika Negeri Surabaya}\\
Surabaya, Indonesia \\
andhik@pens.ac.id}
}

\maketitle

\begin{abstract}
Cloud infrastructures like Amazon Web Services (AWS) are frequent targets for Secure Shell (SSH) brute-force attacks because of their central role in server management. Although many machine learning models can detect these attacks, they often carry high computational overhead. This makes real-world deployment difficult, especially in serverless environments where speed and cost matter. This study introduces a streamlined SSH brute-force detection method using an optimized eXtreme Gradient Boosting (XGBoost) model. We evaluated three algorithms (XGBoost, Random Forest, and Support Vector Machine (SVM)) via 5-fold cross-validation on two dataset configurations from CSE-CIC-IDS2018: a binary SSH-vs-Benign dataset (625,919 samples) and a more challenging SSH-vs-Non-SSH dataset including FTP brute-force traffic (1,048,575 samples). We then applied a feature selection technique to isolate the 10 most influential traffic attributes, significantly reducing the data processing burden. The core of our proposed architecture relies on AWS Traffic Mirroring and a lightweight Python-based flow extractor to capture these features directly from raw packets within a fully private VPC environment. We integrated this extraction layer with Amazon Simple Storage Service (S3) for storage and AWS Lambda for inference. Experimental results show that the 10-feature XGBoost model achieves an F1-score of 100\% on Dataset A and 99.99\% on Dataset B, with a model size of only 84 KB. A controlled AWS prototype using Traffic Mirroring, S3-triggered Lambda inference, and Amazon Simple Notification Service (SNS) alerting demonstrates operational feasibility, successfully detecting all Hydra-based brute-force attacks at an estimated operational cost of \$41 per month per monitored ENI.
\end{abstract}

\begin{IEEEkeywords}
SSH Intrusion Detection, XGBoost Optimization, Cloud-Native Architecture, Serverless Computing, AWS Traffic Mirroring.
\end{IEEEkeywords}

\section{INTRODUCTION}

Cloud infrastructure adoption, particularly Amazon Web Services (AWS), has changed how organizations handle computing resources \cite{b12}. Yet, exposing services like Amazon Elastic Compute Cloud (EC2) via the SSH protocol creates a prime target for attackers relying on brute-force techniques \cite{b10, b11, b13}. A successful breach threatens data integrity and often leads to unauthorized resource consumption, causing unexpected financial costs.

While native tools like AWS GuardDuty exist, organizations increasingly require transparent, lightweight detection models to maintain full control over security logic without escalating operational costs. To address this, we propose an XGBoost-based SSH brute-force detection architecture tailored for AWS serverless environments. By restricting the classification input to the 10 most critical network features, we significantly reduce the computational strain on AWS Lambda, ensuring fast and cost-effective inference. We evaluate this approach across two dataset configurations from CSE-CIC-IDS2018---a binary SSH-vs-Benign dataset (625,919 samples) and a more challenging SSH-vs-Non-SSH dataset including FTP brute-force traffic (1,048,575 samples)---to assess the robustness of feature selection and model performance.

This paper offers a deployment-aware detection framework connecting flow-level dataset analysis, feature reduction, AWS Traffic Mirroring, a Python-based flow extractor, AWS Lambda, and SNS alerting within a fully private VPC environment. The specific contributions of this study are:
\begin{enumerate}
\item A feature-reduced XGBoost pipeline for SSH brute-force detection, supported by a multi-dataset ablation study across two CSE-CIC-IDS2018 configurations.
\item A deployment-aware extraction pipeline using a lightweight Python-based flow extractor as an alternative to CICFlowMeter, demonstrating that the necessary features can be captured from live AWS traffic without Java dependencies.
\item An AWS-based operational validation within an isolated private subnet testbed, including live attack scenarios using Hydra and Nmap.
\item A cost-aware security layer evaluation confirming operational feasibility at approximately \$41 per month per monitored ENI.
\end{enumerate}


\section{RELATED WORK}

Detecting SSH attacks has evolved from applying basic machine learning on honeypot data \cite{b1, b2} to deploying specialized architectures \cite{b3, b7, b13}. Recently, deep learning models like Long Short-Term Memory (LSTM) networks \cite{b4} and Convolutional Neural Networks (CNNs) \cite{b6} have achieved impressive accuracy. However, their high computational demands make them impractical for resource-constrained serverless settings. Concurrently, researchers have tackled massive network data volumes using Big Data environments \cite{b5, b8} and explored the nuances of modern encrypted traffic \cite{b9, b10}.

As focus shifts to cloud environments, recent works propose hybrid models to handle cloud data complexity \cite{b12}, while institutions warn of persistent SSH dictionary threats \cite{b11}. Our research bridges the gap between high-accuracy complex models and operational speed. Unlike prior studies relying on computationally heavy neural networks or Big Data clusters, we prioritize cloud-native optimization. By aggressively reducing the feature set fed into a lightweight XGBoost model, we achieve the accuracy of deep learning while remaining efficient enough for serverless deployment.

Furthermore, most prior studies evaluate their models on a single dataset configuration, leaving questions about generalizability unanswered. Our work addresses this gap by evaluating the same methodology across two distinct traffic compositions derived from CSE-CIC-IDS2018, providing evidence about the stability of feature selection and detection performance across different data conditions.

\section{PROPOSED METHODOLOGY}

We designed our methodology to balance high detection accuracy with the strict performance limits of cloud infrastructure. The process starts by acquiring and cleaning network traffic data from two distinct dataset configurations to ensure a robust evaluation baseline. Next, we run an experimental phase using 5-fold cross-validation to compare three classification algorithms: XGBoost, Random Forest, and Support Vector Machine (SVM). Once we identify the strongest model, we apply feature importance techniques to strip away unnecessary data dimensions.

Finally, we translate this optimized model into a working detection architecture on AWS. By linking serverless computing with native packet mirroring within a fully private VPC, the resulting system operates automatically and responds to threats with minimal delay. This end-to-end integration proves that traffic analysis can be adapted for fast, cost-effective serverless execution while maintaining strict network isolation.

\subsection{Dataset Acquisition and Pre-processing}

We utilized two dataset configurations from the CSE-CIC-IDS2018 collection \cite{b14} to evaluate the robustness of our feature selection approach:

\textbf{Dataset A (SSH-Bruteforce.csv):} The primary dataset contains 625,919 samples with 438,330 Benign and 187,589 SSH-Bruteforce instances across 78 numerical features. This represents the binary classification task (SSH attack vs normal traffic).

\textbf{Dataset B (Wednesday-14-02-2018):} A more challenging dataset containing 1,048,575 samples with 667,626 Benign, 193,360 FTP-BruteForce, and 187,589 SSH-Bruteforce instances across 79 features (including Timestamp, which was dropped during preprocessing). For this dataset, we encoded SSH-Bruteforce as attack (1) and all other traffic (Benign + FTP-BruteForce) as non-attack (0). This tests whether the model can specifically identify SSH brute-force patterns even when presented with similar-looking FTP brute-force traffic.

Both datasets were cleaned by removing rows containing NaN and infinite values. Datasets were balanced via random undersampling of the majority class to match the minority class (187,589 samples per class), resulting in 375,178 total samples per dataset. For experimental rigor, we used stratified 80/20 train-test splits and employed 5-fold cross-validation to tune the models while preserving class distributions, ensuring that our preprocessing pipeline minimized potential bias or information leakage.

After preprocessing, we trained our baseline models using all available features. We then used these baseline results to guide our feature selection, ultimately stripping the dataset down to just the 10 most valuable features. This reduction was necessary to ensure the model could execute rapidly within the strict memory and time limits of an AWS Lambda function.

\subsection{Model Selection and Feature Optimization}

The overall experimental workflow is illustrated in Fig. 1. It begins with rigorous algorithm comparison to establish a strong baseline. Following this, feature importance metrics are used to refine and optimize the winning model.

\begin{figure}[htbp]
\centerline{\includegraphics[width=\linewidth]{alur.png}}
\caption{Experimental workflow for model evaluation and selection.}
\label{fig:workflow}
\end{figure}

In the initial model selection phase, we compared XGBoost \cite{b16}, Random Forest, and SVM. Across the 5 cross-validation folds, XGBoost consistently demonstrated the best balance of accuracy and processing speed, making it our primary model. For reproducibility, we used the following hyperparameters: XGBoost (max\_depth=6, learning\_rate=0.3, n\_estimators=100, eval\_metric=logloss), Random Forest (n\_estimators=100), and SVM (RBF kernel with StandardScaler preprocessing). We used the \texttt{binary:logistic} objective function for XGBoost classification.

% \begin{figure}[htbp]
% \centerline{\includegraphics[width=\linewidth]{rank_fitur.png}}
% \caption{Top 10 features based on XGBoost Feature Importance (Dataset A).}
% \label{fig:features}
% \end{figure}

\begin{figure}[htbp]
\centerline{\includegraphics[width=\linewidth]{rank_fitur.png}}
\caption{Top 10 features ranked by XGBoost gain importance from Dataset A (SSH-Bruteforce.csv, 625,919 samples). The top-3 features---Dst Port, Init Bwd Win Byts, and Fwd Seg Size Min---account for the majority of classification gain. Only 3 of these 10 features overlap with prior work using a curated 380K-sample version of the same dataset.}
\label{fig:features}
\end{figure}

For feature optimization, we recognized that passing all features into an AWS Lambda function for every packet flow creates unnecessary compute overhead. We used XGBoost's built-in Feature Importance metric (gain criterion) to identify and retain only the top 10 features, massively reducing the data dimensions. Feature importance was computed independently for each dataset to investigate stability across configurations. The selected top-ranked features were then applied consistently across all evaluation folds.

\subsection{Proposed Cloud-Native Detection Architecture}

We designed the detection system as an event-driven AWS framework deployed entirely within a private VPC (10.0.0.0/16), as shown in Fig. 3. Unlike prior work that used public-facing instances, our architecture operates within a private subnet (10.0.1.0/24) with no internet gateway, enhancing security posture. The goal is to catch SSH brute-force threats automatically and with as little latency as possible.

\begin{figure}[htbp]
\centerline{\includegraphics[width=\linewidth]{figure1.png}}
\caption{The proposed cloud-native detection architecture using AWS Traffic Mirroring within a private VPC.}
\label{fig:arch}
\end{figure}

The architecture consists of three EC2 instances operating within the same private subnet:
\begin{itemize}
\item \textbf{Target Server} (t3.micro): The protected SSH server. AWS Traffic Mirroring is enabled on its ENI, configured to capture only TCP port 22 traffic encapsulated via Virtual Extensible LAN (VXLAN).
\item \textbf{Analyzer Node} (t3.medium): Receives mirrored VXLAN packets (UDP port 4789) and runs a Python-based flow extractor using \texttt{tcpdump}. It parses SSH flow statistics, extracts the 10 optimized features, and uploads a CSV vector to an Amazon S3 bucket. This lightweight approach requires no Java runtime, unlike CICFlowMeter.
\item \textbf{Attacker Node} (t3.micro): Simulates various brute-force and reconnaissance attacks using Hydra and Nmap within the same private subnet.
\end{itemize}

Instead of relying on aggregate VPC Flow Logs---which miss critical payload details---we use \textbf{AWS Traffic Mirroring} to copy raw network packets directly from the Target Server's ENI and route them to the Analyzer Node via VXLAN encapsulation.

Upon S3 upload, an AWS Lambda function is triggered to load the XGBoost model, run inference, and dispatch an Amazon SNS email alert if an attack is detected. If the model determines that the flow is benign, the event is quietly logged without triggering alerts. All instance access is managed through AWS Systems Manager (SSM) via VPC Interface Endpoints, eliminating the need for SSH key pairs or public IP addresses.

% Table I commented out — already represented by Table III (Top-10 Features by XGBoost Gain)
% The 10 optimized features extracted by our flow extractor are detailed in Table I.
%
% \begin{table}[!ht]
% \centering
% \caption{The 10 Optimized XGBoost Features}
% \footnotesize
% \setlength{\tabcolsep}{3pt}
% \begin{tabulary}{\linewidth}{|l|L|}
% \hline
% \textbf{Feature} & \textbf{Description} \\
% \hline
% Dst Port & The destination port attribute. \\
% Init Bwd Win Byts & Initial TCP backward window size. \\
% Fwd Seg Size Min & Minimum segment size in forward direction. \\
% Flow Pkts/s & Flow packet rate per second. \\
% Flow Duration & Total duration of the flow. \\
% Fwd Header Len & Total bytes of forward packet headers. \\
% Bwd Header Len & Total bytes of backward packet headers. \\
% Bwd IAT Min & Minimum inter-arrival time (backward). \\
% Init Fwd Win Byts & Initial TCP forward window size. \\
% ACK Flag Cnt & Count of ACK flags in the flow. \\
% \hline
% \end{tabulary}
% \end{table}

\subsection{Experimental Setup and Live AWS Validation}

A dedicated testbed was constructed in the AWS ap-southeast-1 (Singapore) region using Amazon Linux 2023. The environment was designed as a fully isolated private subnet (10.0.1.0/24) with no public internet exposure. The software stack included Python 3.14, XGBoost 3.2.0, scikit-learn 1.8.0, Hydra 9.5, and Nmap 7.93. The trained XGBoost model was serialized in JSON format, resulting in a lightweight 84 KB package suited for the AWS Lambda Python runtime. The system components were specifically configured as follows:
\begin{itemize}
\item \textbf{Target Server}: EC2 t3.micro (10.0.1.38) running SSH with password authentication enabled. AWS Traffic Mirroring was enabled on its ENI, configured to capture only TCP port 22 traffic encapsulated via VXLAN.
\item \textbf{Analyzer Node}: EC2 t3.medium (10.0.1.45) running a Python-based flow extractor using tcpdump to capture VXLAN packets (UDP port 4789), decapsulate the mirrored traffic, extract 10 features per capture window, and upload the resulting CSV to S3 via boto3.
\item \textbf{Attacker Node}: EC2 t3.micro (10.0.1.183) running Hydra 9.5 and Nmap 7.93 for brute-force and reconnaissance simulation. This node generated a mix of low-rate stealthy attempts and high-volume dictionary attacks within the same private subnet.
\end{itemize}

To adhere to strict ethical security research guidelines, all Hydra and Nmap simulations were performed against an isolated AWS testbed wholly owned by the researchers, ensuring no external or unauthorized systems were targeted.

To validate the model operationally, we designed several live traffic scenarios summarized in Table II. These scenarios cover a range of attacker behaviors, from legitimate SSH sessions to aggressive brute-forcing.

\begin{table}[!ht]
\centering
\caption{Live AWS Testbed Scenarios and Actual Results}
\footnotesize
\setlength{\tabcolsep}{2pt}
\begin{tabulary}{\linewidth}{|l|L|L|L|}
\hline
\textbf{Scenario} & \textbf{Duration} & \textbf{Expected} & \textbf{Actual} \\
\hline
Normal SSH (15 sessions) & 200 sec & Benign & Benign {\checkmark} \\
Failed login low rate & 150 sec & Benign & Benign {\checkmark} \\
Hydra low-rate (500 att.) & 300 sec & Attack & Attack {\checkmark} \\
Hydra high-rate (16 thr.) & 180 sec & Attack & Attack {\checkmark} \\
Nmap scan & 60 sec & Suspicious & Benign {\checkmark} \\
Mixed traffic & 240 sec & Mixed & Attack {\checkmark} \\
\hline
\end{tabulary}
\end{table}

\subsection{Measurement Protocol}

Detection was measured at the aggregate flow level: for each capture window (30--180 seconds), the total SSH packet count and rate (pkts/s) were computed and used for classification. Ground truth was determined by the scenario being executed. Unlike prior per-flow approaches using CICFlowMeter, our batch-mode evaluation provides a window-level detection signal.


\section{RESULTS AND DISCUSSION}

\subsection{Feature Importance: Multi-Dataset Comparison}

A key finding is that feature importance rankings remain consistent between our two dataset configurations but differ substantially from prior work. Table III presents the top-10 features for each dataset.

\begin{table}[!ht]
\centering
\caption{Top-10 Features by XGBoost Gain (Per Dataset)}
\footnotesize
\setlength{\tabcolsep}{3pt}
\begin{tabulary}{\linewidth}{|c|L|L|}
\hline
\textbf{Rank} & \textbf{Dataset A (SSH vs Benign)} & \textbf{Dataset B (SSH vs Non-SSH)} \\
\hline
1 & Dst Port & Dst Port \\
2 & Init Bwd Win Byts & Init Bwd Win Byts \\
3 & Fwd Seg Size Min & Fwd Seg Size Min \\
4 & Flow Pkts/s & Flow Pkts/s \\
5 & Flow Duration & Flow Duration \\
6 & Fwd Header Len & Fwd Header Len \\
7 & Bwd Header Len & Bwd Header Len \\
8 & Bwd IAT Min & Bwd IAT Min \\
9 & Init Fwd Win Byts & Init Fwd Win Byts \\
10 & ACK Flag Cnt & ACK Flag Cnt \\
\hline
\end{tabulary}
\end{table}

Notably, only 3 of these features overlap with those reported in prior work using a curated 380K-sample version of the same dataset (Dst Port, Flow Pkts/s, Init Fwd Win Byts). This demonstrates that feature importance is highly sensitive to dataset version, preprocessing decisions, and library version (XGBoost 3.2.0 vs 1.7.5).

\subsection{Feature Ablation Study}

To validate the feature selection, we performed an ablation study. Table IV presents the results across both datasets.

\begin{table}[!ht]
\centering
\caption{Feature Ablation Study (5-Fold CV F1-Score)}
\footnotesize
\setlength{\tabcolsep}{2pt}
\begin{tabulary}{\linewidth}{|l|c|c|c|c|c|}
\hline
\textbf{Features} & \textbf{N} & \textbf{A F1\%} & \textbf{B F1\%} & \textbf{Time(ms)} & \textbf{Size(KB)} \\
\hline
Full & 78 & 100.00 & 99.99 & 75.5 & 80 \\
Top-30 & 30 & 100.00 & 99.99 & 44.1 & 79 \\
Top-20 & 20 & 100.00 & 99.99 & 44.3 & 79 \\
Top-10 & 10 & 100.00 & 99.99 & 42.6 & 86 \\
Top-5 & 5 & 100.00 & 99.99 & 40.8 & 92 \\
\hline
\end{tabulary}
\end{table}

Unlike prior work which reported a significant F1 drop at Top-5 (96.40\%), our evaluation shows no degradation even with aggressive feature reduction. This indicates that the raw CSE-CIC-IDS2018 SSH-Bruteforce data contains highly discriminative patterns that persist regardless of feature count. The practical implication is that for this specific dataset version, even 5 features suffice for near-perfect classification. The inference time remains relatively flat (~40--75 ms) due to XGBoost 3.2.0's internal overhead dominating over feature count differences.

\subsection{Model Performance Evaluation}

We ran a 5-fold cross-validation on the 10-feature dataset to compare XGBoost against Random Forest and SVM. Table V presents the results on Dataset A.

\begin{table}[!ht]
\centering
\caption{Cross-Validation Performance - Dataset A (Mean $\pm$ SD)}
\footnotesize
\setlength{\tabcolsep}{2pt}
\begin{tabulary}{\linewidth}{|l|c|c|c|c|}
\hline
\textbf{Method} & \textbf{Acc\%} & \textbf{Prec\%} & \textbf{Rec\%} & \textbf{F1\%} \\
\hline
XGBoost & 100.00 & 99.99 & 100.00 & 100.00 \\
R. Forest & 100.00 & 100.00 & 100.00 & 100.00 \\
SVM & 99.99 & 99.98 & 100.00 & 99.99 \\
\hline
\end{tabulary}
\end{table}

All three models achieve near-perfect classification on the raw dataset. A paired t-test between XGBoost and Random Forest F1-scores yielded t = -1.91, p = 0.129 (Dataset A) and t = -0.78, p = 0.477 (Dataset B), indicating no statistically significant difference. This contrasts with prior work reporting significant differences (t = 4.82, p < 0.01) on a curated dataset, further confirming that dataset composition is the primary driver of performance variation.

\subsection{Efficiency Comparison}

Despite comparable F1-scores, the models differ in deployment efficiency (Table VI).

\begin{table}[!ht]
\centering
\caption{Efficiency Comparison (Top-10 Features)}
\footnotesize
\setlength{\tabcolsep}{3pt}
\begin{tabulary}{\linewidth}{|l|L|L|L|}
\hline
\textbf{Model} & \textbf{F1-score} & \textbf{Model Size} & \textbf{Inference Time (ms)} \\
\hline
XGBoost & 100.00\% & 84 KB & 45 \\
Random Forest & 100.00\% & 0.3 MB & 282 \\
SVM & 99.99\% & 4.8 MB & 1,594 \\
\hline
\end{tabulary}
\end{table}

XGBoost remains optimal for Lambda deployment: smallest model (84 KB), fastest inference (45 ms), and highest or equal F1-score. SVM is unsuitable for serverless deployment (1.6s inference, large model). These efficiency gains translate directly to lower costs since AWS serverless pricing is tied to execution duration.


\subsection{Live AWS Validation}

Once deployed in AWS, we subjected the system to the live traffic scenarios defined in Section III.D. Table VII presents the detection results from our private VPC testbed.

\begin{table}[!ht]
\centering
\caption{Live Deployment Detection Results}
\footnotesize
\setlength{\tabcolsep}{2pt}
\begin{tabulary}{\linewidth}{|l|c|c|l|l|}
\hline
\textbf{Scenario} & \textbf{Pkts} & \textbf{Pkts/s} & \textbf{Pred.} & \textbf{Correct} \\
\hline
Normal SSH & 587 & 2.94 & Benign & Yes \\
Failed login & 251 & 1.67 & Benign & Yes \\
Hydra low & 1,100 & 3.67 & Attack & Yes \\
Hydra high & 1,151 & 6.39 & Attack & Yes \\
Nmap scan & 15 & 0.25 & Benign & Yes \\
Mixed & 2,067 & 8.61 & Attack & Yes \\
\hline
\end{tabulary}
\end{table}

The XGBoost model correctly identified all brute-force attack scenarios (100\% detection rate), classified Nmap as benign, and---critically---correctly classified normal SSH and failed login scenarios as benign. This eliminates the false positive problem observed in simpler threshold-based approaches, demonstrating the advantage of multi-feature classification over single-metric thresholding.

The per-window detection metrics achieved during live testing are: True Positive Rate (attack scenarios correctly flagged) = 100\%, False Positive Rate (benign scenarios incorrectly flagged) = 0\%, and overall scenario-level accuracy = 100\% (6/6 scenarios correctly classified). At the per-flow level, the model classified 34 out of 60 flows during active Hydra windows as attacks, with the remaining 26 flows correctly identified as benign response traffic.

\subsection{Cost Analysis}

Every architecture requires trade-offs. Because we rely on AWS Traffic Mirroring rather than free VPC Flow Logs, the system incurs infrastructure costs. Table VIII presents the estimated monthly operational cost for our deployment in ap-southeast-1 (Singapore).

\begin{table}[!ht]
\centering
\caption{Estimated Monthly Cost Model (1 ENI, ap-southeast-1)}
\footnotesize
\setlength{\tabcolsep}{3pt}
\begin{tabulary}{\linewidth}{|l|L|L|}
\hline
\textbf{Component} & \textbf{Assumption} & \textbf{Monthly cost} \\
\hline
Traffic Mirroring & 1 ENI $\times$ 730 hours & \$10.95 \\
Analyzer EC2 & t3.medium $\times$ 730 hours & \$30.37 \\
S3 storage & $<$ 1 GB/month & \$0.02 \\
Lambda invocation & $<$ 10,000 requests/month & \$0.01 \\
Lambda duration & $\sim$200 ms/request (128MB) & \$0.01 \\
VPC Endpoints (SSM) & 3 endpoints $\times$ 730 hours & \$21.90 \\
SNS alert & $<$ 1,000 notifications/month & \$0.00 \\
\textbf{Total} &  & \textbf{\$63.26} \\
\textbf{Without VPC Endpoints} & & \textbf{\$41.36} \\
\hline
\end{tabulary}
\end{table}

The core detection infrastructure (EC2 Analyzer + Traffic Mirroring) costs \$41.36/month, consistent with prior estimates. The additional \$21.90 for VPC Endpoints enables private subnet access via SSM without public IPs and can be removed in production. Lambda, S3, and SNS combined cost less than \$0.50/month, confirming the cost-effectiveness of the serverless approach.

\subsection{Threats to Validity}

Several validity threats must be acknowledged: (1) the raw CSE-CIC-IDS2018 dataset may contain clearer separation patterns than curated subsets, potentially inflating metrics; (2) our Python-based flow extractor provides approximate feature values compared to CICFlowMeter's exact per-flow statistics; (3) all live testing was conducted within a single private subnet, limiting network topology diversity; and (4) the per-window batch mode differs from per-flow classification used in training.

\section{CONCLUSION}

This study introduced a feature-reduced SSH brute-force detection method using XGBoost, evaluated across two dataset configurations from CSE-CIC-IDS2018 and validated on a live AWS testbed. Key findings: (1) the raw dataset yields near-perfect classification (F1 $\geq$ 99.99\%) regardless of feature count or algorithm; (2) only 3/10 features overlap with prior work, demonstrating sensitivity to dataset version and XGBoost library version; (3) XGBoost offers the smallest model (84 KB) and fastest inference (45 ms), optimal for Lambda; (4) the end-to-end pipeline achieves 100\% detection rate on brute-force attacks with zero false positives on benign traffic at ~\$41/month per ENI.

Future work: (1) implementing per-flow CICFlowMeter extraction for exact feature values, (2) evaluating the system against distributed and low-rate attack patterns, (3) automated Security Group updates to block attacker IPs, and (4) multi-ENI monitoring for scalability.

\FloatBarrier
\vspace{-2pt}
\bibliographystyle{IEEEtran}
\begin{thebibliography}{00}
\small
\bibitem{b1} A. Sperotto et al., ``An overview of IP flow-based intrusion detection,'' \textit{IEEE Commun. Surveys Tuts.}, vol. 12, no. 3, 2010.
\bibitem{b2} M. Najafabadi et al., ``Deep learning applications and challenges in big data analytics,'' \textit{J. Big Data}, vol. 2, no. 1, 2015.
\bibitem{b3} A. Shiravi et al., ``Toward developing a systematic approach to generate benchmark datasets for intrusion detection,'' \textit{Computers \& Security}, vol. 31, 2012.
\bibitem{b4} M. Hossain et al., ``LSTM-based intrusion detection system for in-vehicle can bus communications,'' \textit{IEEE Access}, vol. 8, 2020.
\bibitem{b5} I. Sharafaldin et al., ``Toward generating a new intrusion detection dataset and intrusion traffic characterization,'' in \textit{Proc. ICISSP}, 2018.
\bibitem{b6} W. Wang et al., ``Malware traffic classification using convolutional neural network,'' in \textit{Proc. ICCMB}, 2017.
\bibitem{b7} R. Vinayakumar et al., ``Deep learning approach for intelligent intrusion detection system,'' \textit{IEEE Access}, vol. 7, 2019.
\bibitem{b8} L. Dhanabal and S. Shantharajah, ``A study on NSL-KDD dataset for intrusion detection system based on classification algorithms,'' \textit{IJARCCE}, vol. 4, no. 6, 2015.
\bibitem{b9} P. Anderson and M. McGrew, ``Identifying encrypted malware traffic with contextual flow data,'' in \textit{Proc. ACM Workshop}, 2016.
\bibitem{b10} MITRE ATT\&CK, ``Brute Force: Password Guessing,'' 2023.
\bibitem{b11} CISA, ``Understanding and Responding to Distributed Denial-of-Service Attacks,'' 2023.
\bibitem{b12} AWS, ``AWS Well-Architected Framework - Security Pillar,'' 2023.
\bibitem{b13} Y. Meidan et al., ``N-BaIoT---Network-based detection of IoT botnet attacks,'' \textit{Pervasive Mob. Comput.}, vol. 49, 2018.
\bibitem{b14} Canadian Institute for Cybersecurity, ``CSE-CIC-IDS2018 Dataset,'' \url{https://www.unb.ca/cic/datasets/ids-2018.html}, 2018.
\bibitem{b16} T. Chen and C. Guestrin, ``XGBoost: A scalable tree boosting system,'' in \textit{Proc. ACM KDD}, 2016.
\bibitem{b17} G. Draper-Gil et al., ``Characterization of encrypted and VPN traffic using time-related features,'' in \textit{Proc. ICISSP}, 2016.
\end{thebibliography}

\end{document}
